\section{Discussion}

NeuRev Workbench treats scientific interpretation as a provenance problem as
well as a modeling problem. In the present recording, the same pipeline can
produce a recovered known-positive occurrence, an unmatched candidate for expert
review, a ranking miss, a proposal-absence case, or an apparent recentering gain
that actually crosses identity. Representing all of these outputs as detections
would discard the information needed to audit or improve the system.

\subsection{Identity is a first-class data structure}

The empirical recentering analysis shows why observation site, geometry, and
canonical identity should not be stored as one mutable record.
\VenueIdentityCollision{} of \VenueStrictMissed{} local improvements overlapped
an existing identity. A software system that updates the
center in place could convert those cases into apparent recoveries without
leaving evidence of reassignment. Immutable observations and explicit identity
relationships instead permit alternative geometries to be evaluated while the
original reference remains recoverable.

This pattern extends beyond calcium imaging. Any workflow with crowded objects,
repeated observations, uncertain matching, or human adjudication benefits from
separating what was measured from what it is believed to represent. Geometry
hashes and versioned identity decisions provide a practical basis for that
separation.

\subsection{Unknown is an operational state, not missing metadata}

Positive--unlabeled evaluation is often introduced as a statistical choice. Here
it is also enforced as a data-state contract. Unmatched candidates stay unknown
through feature computation, out-of-fold scoring, review export, and manuscript
generation. This prevents downstream tables from silently converting missing
labels into negatives.

The \VenueReviewedCandidates{}-site review illustrates the benefit.
\VenueReviewLikely{} calls were definite or probable, but their selected-batch
frequency was never routed into a precision field. \VenueReviewUncertain{}
uncertain calls remained unresolved rather than negative. The
priority score was retained as review value despite poor alignment with obvious
likely-neuron calls, because its intended target was ambiguity and crowding.
Explicit task semantics make that result interpretable rather than erroneous.

\subsection{Evidence-bound generation limits narrative drift}

The workflow separates run completion, artifact validation, scientific-audit
completion, and scientific promotion. This distinction was consequential during
manuscript audit: a completed local Feature Atlas run had numerical artifacts but
no canonical experiment/run binding and an incomplete model-annotation media
audit. The corrected build excludes its values from both abstracts, main Results,
and main figures. Completion is visible as engineering history without being
misrepresented as a registered model replacement.

Generating synchronized technical, overview, and experiment-history views from
the same registry reduces copy-and-paste drift. The venue layer adds a compact
source-to-fact manifest because not every maintained result is yet represented
by a portable registry capsule. Quantitative facts, scope, evidence links, and
status are appropriate generated fields; nuanced interpretation remains in
author-reviewed prose.

\subsection{Software direction}

After independent adjudication and a bounded truth set, a useful next software
component to test is a competition-aware identity resolver.
Rather than producing one scalar score per location, it should construct local
candidate clusters and represent retain, merge, split, or abstain decisions.
Cluster size, footprint overlap, peak separation, trace partial correlation,
within-cluster score entropy, and assignment stability are candidate inputs.
Evaluation should preserve identity-level recovery, duplicate count, and review
coverage as separate outputs.

Future resolver studies should include quality-control or abstention fields and
keep their intended role separate from neuron-probability ranking. Future
releases should also provide portable example data, a documented command-line
workflow, machine-readable schemas, deterministic smoke tests, and archive
identifiers for code and artifacts.

\subsection{Limitations and generalization}

The current demonstration uses one recording and historically informed event
windows. It does not establish cross-recording transfer or broad software utility.
The candidate review has one expert and awaits adjudication. Full-field precision
requires a bounded exhaustive truth set. Some analyses have few independent
identity groups, while exact-truth simulations do not establish biological
transfer. The public software/data release and clean release-suite
record are also pending.

These limitations define the next validation sequence: freeze the identity and
candidate contracts, complete independent review, construct bounded truth without
exposing model scores during raw-first annotation, and then transfer the frozen
workflow to another recording. A successful result would support a general
neuroinformatics contribution; a failed transfer would still localize which
data, acquisition, or identity assumptions did not generalize.
